\documentclass[11pt]{amsart}
\usepackage[T1]{fontenc}
\usepackage[utf8]{inputenc}
\usepackage{amssymb}
\usepackage[margin=1.05in]{geometry}
\usepackage[hidelinks]{hyperref}
\newcommand{\stamptime}{\the\numexpr\time/60\relax:\ifnum\numexpr\time-60*(\time/60)\relax<10 0\fi\the\numexpr\time-60*(\time/60)\relax}
\AddToHook{shipout/foreground}{\put(0.5\paperwidth,-\paperheight+0.55cm){\makebox[0pt]{\tiny\ttfamily unit/doc/\jobname.tex \ --- \ compiled \today, \stamptime}}}
\newtheorem{theorem}{Theorem}
\newtheorem{proposition}[theorem]{Proposition}
\newtheorem{lemma}[theorem]{Lemma}
\newtheorem{corollary}[theorem]{Corollary}
\theoremstyle{definition}
\newtheorem{definition}[theorem]{Definition}
\newtheorem{construction}[theorem]{Construction}
\theoremstyle{remark}
\newtheorem{remark}[theorem]{Remark}
\newcommand{\Om}{\Omega}
\newcommand{\Map}{\mathrm{Map}_*}
\newcommand{\Z}{\mathbb{Z}}
\newcommand{\code}[1]{\texttt{#1}}
\newcommand{\cT}{\mathsf{isTrunc}}
\newcommand{\cC}{\mathsf{isConn}}
\newcommand{\cc}{\cdot}
\newcommand{\refl}{\mathsf{refl}}
\newcommand{\ev}{\mathsf{ev}}
\newcommand{\lp}{\mathsf{loop}}
\title{The cup square $\kappa:K(\Z,2)\to_*K(\Z,4)$ by delooping}
\date{28 October 2026}
\begin{document}
\maketitle

\noindent\textbf{Purpose and status.} Informal construction for step
\S2b of \code{doc/pi3-s2-postnikov-handoff-3.txt}: the pointed map
$\kappa$ that \code{KZ2Cohomology.H4} is about, built as the cup square
$x\mapsto x\smile x$. Written before the formal proof; each numbered
statement is a declaration of \code{CupSquare}
(\S\ref{sec:corr}). Proving (H4) for this $\kappa$ is the next step (\S2c)
and is not attempted here. Foundations: those of the development
(univalence, images and their surjectivity; function extensionality
derived). No new axiom; the sphere axiom is not used.
Milestone: \emph{a constructed, normalised cup square $\kappa$}; that
coefficient change along $\kappa$ is a bijection is the separate theorem
(H4). \emph{Status: checked} (64 modules, 0 failed; 2380 declarations;
\code{cup.opens}: nothing else opened). Refereed independently (fresh reader given the note and
the library, not the conclusion): constructions correct, interfaces and
indices confirmed; corrected here: the centre of $S$ is not definitional
(recentring, \code{contr-at}), lemmas added to (I6), the direction
$f=\lp\,c_2$, and the remark on (H4) (the looping of $\kappa$ is null).

The construction is that of Wärn (\emph{Eilenberg--MacLane spaces and
stabilisation in homotopy type theory}, J.\ Homotopy Relat.\ Struct.\ 2023,
\S4.3): the cup product $K(A,n)\to_*\Map(K(B,m),K(C,n+m))$ is obtained by
delooping, one level at a time, starting from a bilinear map
$A\times B\to C$. We need only $A=B=C=\Z$, multiplication, and the two
levels $n=1,2$ with $m=2$; for $A=\Z$ the first level is given by the
circle property of $K(\Z,1)$ (\code{CircleProperty}, step \S2a), so no
coefficient maps $K(a\cdot{-},2)$ and no multiplicativity argument are
needed.

\section{Conventions and interfaces}

As in \code{doc/pi3-s2-postnikov.pdf}: \code{isTrunc n X} is ``$X$ is
$(n-2)$-truncated'', \code{isConn n X} is ``$X$ is $(n-2)$-connected'';
$K_m=K(\Z,m)$ is \code{EMMaps.KK[Zab].K m}, pointed by $*$
(\code{KK.pt m}); $\Om Y=(y_0=y_0,\refl)$. Paths compose in diagrammatic
order, $p\cc q$ is \code{trans p q}, by $J$ on $q$, so $p\cc\refl\equiv p$.
For pointed $(X,x_0)$, $(Y,y_0)$ the pointed mapping type
$\Map(X,Y)=(F:X\to Y)\times(F\,x_0=y_0)$ (\code{Lemma4.PtMap}) is itself
pointed, by the constant map $0=(\lambda\_.\,y_0,\refl)$.

The construction uses the following, all in the library or proved in \S2a:
\begin{enumerate}
\item[(I1)] \code{KK}: $\cC\,(m+2)\,K_{m+1}$, $\cT\,(m+3)\,K_{m+1}$,
  and pointed equivalences $e_m:\Om K_{m+1}\simeq_*K_m$ (\code{KK.e}).
\item[(I2)] \code{Lemma4.cor5-map}: if $\cC\,(c+1)\,X$ and
  $\cT\,(c+t)\,Y$ then $\cT\,t\,\Map(X,Y)$.
\item[(I3)] \code{DeloopMaps.Maps.Over.cor11} (DW1 Cor.~11): if
  $\cC\,(c+2)\,X$ and $\cT\,(2(c+1))\,Y$, then for every pointed
  $f:\Om X\to_*\Om Y$ the type
  $\mathsf{Unloop}\,f=(F:\Map(X,Y))\times(f=\lp\,F)$ is contractible.
\item[(I4)] \S2a, \code{CircleProperty.Circle}: for pointed $(Y,y_0)$,
  $\ev_\circ:\Map(K_1,Y)\simeq\Om Y$, $F\mapsto(\lp\,F)\,g$, where $g$ is
  the generator of $\Om K_1$ (\code{pointed-equiv}).
\item[(I5)] \code{LoopIso.loopPE} ($\Om$ of a pointed equivalence),
  \code{PtEquivOps.pt-comp}, \code{pt-inv}; \code{LoopCalc.ptcomp}.
\item[(I6)] \code{FundThm.FT} (fundamental theorem of identity types;
  the map is transport, so it sends $\refl$ to the given element by
  definition), \code{FunExtUA.FE.wfe} (a product of contractible types is
  contractible), the singletons \code{DeloopMaps.singL} ($(y:Y)\times(y_0=y)$)
  and \code{FunExtUA.sing} ($(w:W)\times(w=w_0)$), \code{FunExtUA.pathC},
  \code{Fibres.trans-equiv} ($a:x=y$ gives $(y=z)\simeq(x=z)$, $q\mapsto a\cc q$),
  \code{HLevels.trunc-Sigma-contr} (the fibre over the \emph{centre carried
  by the contractibility proof} decides), \code{Retracts} (\code{rt-Sigma},
  \code{rt-trunc}) and \code{CircleProperty.rt-equiv} (an equivalence as a
  retract).
\end{enumerate}
The one new general lemma is Lemma~\ref{lem:loops}.

\section{Loops of a pointed mapping type}

\begin{lemma}\label{lem:loops}
For pointed $(X,x_0)$, $(Y,y_0)$, with $M=\Map(X,Y)$ pointed by $0$,
there is a pointed equivalence
\[\ev:\Om M\simeq_*\Map(X,\Om Y),\qquad \ev\,\refl\equiv r
  =(\lambda\_.\,\refl_{y_0},\refl),\]
the base point of $\Map(X,\Om Y)$, which evaluates a path of pointed maps
pointwise: for $P:0=G$ and $x:X$, $(\ev\,P)_1\,x=\mathsf{ap}_{\lambda F.\,F_1x}\,P$
(\code{ev-ap}, by $J$ on $P$: both sides are $\refl$ at $\refl$).
\end{lemma}
\begin{proof}
For $G=(k,w):M$ let
\[R\,G=\bigl(H:(x:X)\to y_0=k\,x\bigr)\times\bigl(H\,x_0\cc w=\refl_{y_0}\bigr).\]
At $G=0$ we have $w=\refl$, so $H\,x_0\cc w\equiv H\,x_0$ and $R\,0$ is,
by definition, $\Map(X,\Om Y)$; let $r=(\lambda\_.\refl,\refl):R\,0$, its
base point. By the fundamental theorem of identity types (I6), transport
$\mathsf{tr}_G:(0=G)\to R\,G$, $p\mapsto p_*r$, is an equivalence as soon
as the total space $T=(G:M)\times R\,G$ is contractible; take
$\ev=\mathsf{tr}_0$. Then $\ev\,\refl\equiv r$, so $\ev$ is pointed by
$\refl$, and it is an equivalence.

$T$ is contractible. Reordering components (a retract whose round trip is
the identity by definition),
\[T\ \lhd\ (s:S)\times W\,s,\quad S=(k:X\to Y)\times((x:X)\to y_0=k\,x),\quad
W(k,H)=(w:k\,x_0=y_0)\times(H\,x_0\cc w=\refl).\]
$S$ is a retract (again by definition) of $(x:X)\to(y:Y)\times(y_0=y)$,
which is contractible by \code{wfe} and \code{singL}. The centre produced
this way goes through univalence and is not $(\lambda\_.\,y_0,\lambda\_.\,\refl)$
by definition, so we \emph{recentre}: a contractible type is contractible
with any given centre $a$, with contraction $z\mapsto\mathsf{pathC}\,a\,z$
(\code{contr-at}). With centre $s_0=(\lambda\_.\,y_0,\lambda\_.\,\refl)$ the
fibre is, by definition, $W\,s_0=W_0=(w:y_0=y_0)\times(\refl\cc w=\refl)$.
For each $w$, composition on the left with
$a_w=(\mathsf{trans\text{-}refl\text{-}l}\,w)^{-1}:w=\refl\cc w$ is an
equivalence $(\refl\cc w=\refl)\simeq(w=\refl)$ (\code{trans-equiv}), hence
a retract (\code{rt-equiv}; its round trip is a path, not a definition);
by \code{rt-Sigma}, $W_0$ is a retract of the singleton
$(w:y_0=y_0)\times(w=\refl)$ (\code{sing}), contractible. Hence
$(s:S)\times W\,s$ is contractible (\code{trunc-Sigma-contr}), and so is
$T$ (\code{rt-trunc}).
\end{proof}

\begin{remark}
Only the pointed inverse $\ev^{-1}:\Map(X,\Om Y)\to_*\Om M$
(\code{pt-inv}) is used in the construction; the equivalence itself
(and the formula $\ev\,P=P_*r$) is what (H4) will use to compute loops of
$\Map(K_2,K_4)$.
\end{remark}

\section{The construction}

Fix $X=K_2$, pointed by $x_0=*$, and write $M(Y)=\Map(K_2,Y)$, pointed by $0$.

\begin{construction}[the fundamental class]\label{c:iota}
$\iota:K_2\to_*\Om^2K_4$ is the inverse (\code{pt-inv}) of the pointed
equivalence
\[\Om^2K_4\ \xrightarrow{\ \Om e_3\ }\ \Om K_3\ \xrightarrow{\ e_2\ }\ K_2\]
(\code{loopPE}, \code{pt-comp}).
\end{construction}

\begin{construction}[level one]\label{c:one}
Let $\ell=\ev^{-1}\,\iota:\Om M(\Om K_4)$ (Lemma~\ref{lem:loops} for
$Y=\Om K_4$), and let
\[c_1=\ev_\circ^{-1}\,\ell\ :\ K_1\to_*M(\Om K_4)\]
by the circle property (I4) for $Y=M(\Om K_4)$, $y_0=0$. By the section
equation, $(\lp\,c_1)\,g=\ell$.
\end{construction}

\begin{construction}[level two: delooping]\label{c:two}
Let $f:\Om K_2\to_*\Om M(K_4)$ be the composite of pointed maps
\[\Om K_2\xrightarrow{\ e_1\ }K_1\xrightarrow{\ c_1\ }M(\Om K_4)
  \xrightarrow{\ \ev^{-1}\ }\Om M(K_4)\]
(Lemma~\ref{lem:loops} for $Y=K_4$). Apply (I3) with $c=1$:
$K_2$ is 1-connected, $\cC\,3\,K_2$ (I1), and $M(K_4)$ is 2-truncated,
$\cT\,4\,M(K_4)$, by (I2) with $c=2$, $t=4$, since
$\cT\,6\,K_4$ (I1). So $\mathsf{Unloop}\,f$ is contractible; let
\[c_2:K_2\to_*M(K_4),\qquad f=\lp\,c_2\]
be its centre (both components; the direction is that of \code{Unloop}).
Informally, $c_2$ is $x\mapsto(x\smile{-})$; nothing below uses this gloss.
\end{construction}

\begin{definition}[the cup square]\label{d:kappa}
$\kappa:K_2\to_*K_4$, $\kappa\,x=(c_2\,x)_1\,x$, pointed by the pointing
$(c_2\,x_0)_2:(c_2\,x_0)_1\,x_0=*$ of the map $c_2\,x_0$.
\end{definition}

\begin{theorem}[normalisation of $c_2$]\label{t:norm}
Let $\bar c_1:K_1\to_*M(\Om K_4)$ be $c_2$ looped and read through the
equivalences: $\bar c_1=e_1^{-1};\lp\,c_2;\ev$ (\code{pt-inv}, \code{ptcomp},
Lemma~\ref{lem:loops} for $Y=K_4$). Then $\ev\,(\ev_\circ\,\bar c_1)=\iota$
(Lemma~\ref{lem:loops} for $Y=\Om K_4$): looping $c_2$, applying the
mapping-loop equivalences and evaluating on the circle generator $g$ gives
back the fundamental class. This records which cup product is constructed.
\end{theorem}
\begin{proof}
The underlying maps are homotopic: for $k:K_1$, using $f=\lp\,c_2$,
$\ev\circ\ev^{-1}=\mathrm{id}$ and $e_1\circ e_1^{-1}=\mathrm{id}$ (section
equations),
$(\bar c_1)_1\,k=\ev((\lp\,c_2)_1(e_1^{-1}k))=\ev(f_1(e_1^{-1}k))
=\ev(\ev^{-1}(c_{1}(e_1(e_1^{-1}k))))=c_1\,k$. We do not claim a pointed
equality: its coherence would be an equation between two paths in
$M(\Om K_4)$, which is a 1-type, not a set. Instead, for pointed maps with
an unpointed homotopy $H$, the loopings are conjugate:
$\lp\,\bar c_1\,q=c^{-1}\cc\lp\,c_1\,q\cc c$ for a loop $c$ of $M(\Om K_4)$
at $0$ (\code{LoopConj.LoopHomotopy.loop-homotopy}). The type $M(\Om K_4)$
is pointed equivalent to the loop space $\Om M(K_4)$
(Lemma~\ref{lem:loops} for $Y=K_4$), so its loops commute as paths
(Eckmann--Hilton, \code{EH.Abelian.comm}) and conjugation is the identity.
Hence $\ev_\circ\,\bar c_1=(\lp\,\bar c_1)\,g=(\lp\,c_1)\,g=\ell$
(\code{evPt-loop}), and $\ev\,\ell=\iota$ by the section equation.
\end{proof}

\begin{remark}[canonicity and relation to Wärn]\label{r:canon}
Every choice above is forced by $\iota$ and the generator
$g=\mathsf{unband}\,1$ of $\Om K_1$ used in (I4): $\ev$ and $\ev_\circ$ are
equivalences and $\mathsf{Unloop}\,f$ is contractible, so $c_2$ is the
unique pointed map whose looping is $f$, and $\kappa$ depends only on the
pointed equivalences $e_m$ of (I1) and on $g$. In Wärn's construction for the
bilinear map $\Z\times\Z\to\Z$, level zero is $a\mapsto K(a\cdot{-},2)$,
a homomorphism $\Z\to\Om M(\Om K_4)$ (after $e$); for the group $\Z$ a
homomorphism is given by the image of $1$, $K(1\cdot{-},2)=\mathrm{id}$,
which corresponds to $\ell$; the circle property replaces the delooping
$K(\Z,0)\to K(\Z,1)$ (DW1 Cor.~12 and its multiplicativity condition).
Level two is Wärn's delooping step, by (I3). The pointing of $\kappa$
could also be taken from the pointing $c_2\,x_0=0$ of $c_2$, evaluated at
$x_0$; the two pointings differ by a path, which (H4) may have to
relate. We fix the pointing of the map $c_2\,x_0$, a component, which
needs no lemma.
\end{remark}

\begin{remark}[what (H4) will need; the looping of $\kappa$ is null]
The characterising equations of the construction are
$\ev\,\refl\equiv r$, $(\lp\,c_1)\,g=\ell$ (section equation, then
\code{evPt-loop}) and $f=\lp\,c_2$ (I3). They describe $c_2$, not
$\kappa$ directly: the looping of $\kappa$ is
$\lp\,\kappa\,p=(\mathsf{ap}_{c_2}p)(x_0)\cc\mathsf{ap}_{c_2x_0}p$, whose
first factor is trivial because the points of $M(K_4)$ are pointed maps
and $c_2\,x_0=0$. Indeed $\Om\kappa:K_1\to_*K_3$ is null
($[K(\Z,1),K(\Z,3)]_*=H^3(S^1)=0$). So looping cannot detect the class of
$\kappa$. (H4) must identify $[\kappa]$ in $[K_2,K_4]_*$ (the set of
deloopings of the zero map $K_1\to_*K_3$, \code{doc/pi3-s2-postnikov.pdf})
using $c_2$ as a map $K_2\to_*M(K_4)$, i.e.\ its effect on $\pi_2$ (by the
equations above, through $c_1$ and $\iota$), together with the diagonal;
and it must prove naturality of $g\mapsto[\kappa;K(g,4)]$ in the
coefficients, which this construction does not provide.
\end{remark}

\section{Correspondence with the formal proof}\label{sec:corr}
File \code{CupSquare.seg} (after \code{CircleProperty} in \code{All.seg}).
File \code{sag/pi3-s2-postnikov/CupSquare.seg}; \code{--opens} specification
\code{cup.opens}.
\begin{itemize}
\item Lemma~\ref{lem:loops}: \code{contr-at}; module \code{PtMapLoops[X,x0,Y,y0]}:
  \code{pt0}, \code{R}, \code{r}, \code{R-pt0} (by \code{refl}), \code{S},
  \code{s0}, \code{W}, \code{S-rt}, \code{S-contr}, \code{W0-rt}, \code{T-rt},
  \code{T-contr}, \code{ev} (\code{FT.tr pt0}), \code{ev-ap}, \code{ev-pt}
  (pointed equivalence), \code{unev} (its pointed inverse), \code{unev-un},
  \code{ev-unev} (then \code{unev} is locked).
\item Construction~\ref{c:iota}: \code{iota-eq}, \code{iota}.
\item Construction~\ref{c:one}: \code{ell}, \code{ell-ev}, \code{c1},
  \code{c1-ev} (then \code{c1} is locked).
\item Construction~\ref{c:two}: \code{e1} (locked), \code{f}, \code{tM4},
  \code{cell}, \code{c2}, \code{c2-loop} (then \code{c2} is locked).
\item Definition~\ref{d:kappa}: \code{kappa : PtMap KZ2 kz2 KZ4 kz4}.
\item Theorem~\ref{t:norm}: \code{e1inv}, \code{e1-back}, \code{c1bar},
  \code{c1bar-hom}, \code{AB} (\code{EH.Abelian}), \code{conj-triv},
  \code{LH} (\code{LoopConj.LoopHomotopy}), \code{cup-norm}.
\end{itemize}
\end{document}
